\documentclass{article}

% Pass 'final' to the neurips_2023 package to remove line numbers if needed
\usepackage[final]{neurips_2023}

\usepackage{graphicx}

\usepackage{enumitem}

\usepackage[utf8]{inputenc} % allow utf-8 input
\usepackage[T1]{fontenc}    % use 8-bit T1 fonts
\usepackage{hyperref}       % hyperlinks
\usepackage{url}            % simple URL typesetting
\usepackage{booktabs}       % professional-quality tables
\usepackage{amsfonts}       % blackboard math symbols
\usepackage{nicefrac}       % compact symbols for 1/2, etc.
\usepackage{microtype}      % microtypography
\usepackage{xcolor}         % colors

\usepackage{hyperref}
\usepackage{cleveref}

\title{Summarization: Explicitly Encouraging Low Fractional Dimensional Trajectories via Reinforcement Learning}

% Paper BibTex
% @inproceedings{gillen2021explicitly,
%   title={Explicitly encouraging low fractional dimensional trajectories via reinforcement learning},
%   author={Gillen, Sean and Byl, Katie},
%   booktitle={Conference on Robot Learning},
%   pages={2137--2147},
%   year={2021},
%   organization={PMLR}
% }

\author{
  Summarized by Xiaonan (Nice) Wang \thanks{wangxiaonannice@gmail.com} \\
  % Paper fr last-names (year)
  Summarization\thanks{\textbf{Disclaimer:} This summarization is for research and study purposes. It represents a personal interpretation and may contain inaccuracies. Feedback or corrections via email are highly appreciated.} Generated based on Paper from Gillen and Byl. (2021) \\
  Topic: Reinforcement Learning \\
}

\begin{document}

\maketitle

\begin{abstract}
This document summarizes the core contributions and methodology of the paper "Explicitly Encouraging Low Fractional Dimensional Trajectories via Reinforcement Learning, \cite{gillen2021explicitly}", focusing on its' main ideas and the core blocks.
\end{abstract}

\section{Overview: Core Questions and Answers}
                           
\paragraph{(1) What is the problem?} 
\begin{itemize}
  \item \textbf{Nonlinear} and \textbf{hybrid} system for \textbf{humanoid running}.
\end{itemize}

\paragraph{(2) Why need to solve this problem?}
\begin{itemize}
  \item Achieving performant and robust running;
  \item Rejecting disturbances.
\end{itemize}

\paragraph{(3) How is it different from prev.?}
\begin{itemize}
  \item Embed multi-domain trajectories and CLFs directly into the RL reward, \textbf{no} trajectories or CLFs at \textbf{runtime};
  \item Achieve \textbf{certifiable stability}.
\end{itemize}

\paragraph{(4) Why is it better than prev.? (Advantages)} 
\begin{itemize}
  \item Compared with \textit{heuristically designed controller}: can succeed in \textbf{Hybrid Zero Dynamics (HZD)};
  \item Compared with \textit{offline trajectory optimization + online tracking with traditional feedback linearization \& control Lyapunov functions (CLFs)}: operate on both \textbf{continuous} and \textbf{discrete} (i.e. hybrid) dynamics, with capacity to generate \textbf{transient} behaviors;
  \item Compared with \textit{(contact) impact invariance} solutions like \textit{MPC}: yield humanoid running and \textbf{computational efficient};
  \item Compared with \textit{end-to-end RL}: \textbf{eliminate} the need to \textbf{handcraft} and \textbf{tune heuristic} reward terms;
  \item Produce \textbf{steady state} running motions with accurate position and velocity tracking.
\end{itemize}

\paragraph{(5) What is the approach itself?}

\paragraph{(6) What are the applications of it?} 
\begin{itemize}
  \item A \textbf{Hybrid Controller} that can be used both in simulation and real-world \textbf{Humanoid} Robots (eg. UniTree G1, Atlas) for \textbf{periodic steady state} \textit{running}.
\end{itemize}

\section{The Structure}

\paragraph{Summarized Block-Diagram} 
The summarized block-diagram of \textit{CLF-RL controller synthesizing framework} see \cref{fig:clf_rl}\footnote{For efficiency, this diagram was initially hand-drawn on paper and then converted into its current digital version using Nano Banana Pro.}.

\begin{figure}[tb]
  \centering
  \includegraphics[width=0.95\textwidth]{fig/clf_rl.png}
  \caption{Summarized Block-Diagram of CLF-RL Controller Synthesizing Framework}
  \label{fig:clf_rl}
\end{figure}

\paragraph{Original in Paper}

The original block-diagram see \cref{fig:original}.

\begin{figure}[tb]
  \centering
  \includegraphics[width=0.95\textwidth]{fig/orig.png}
  \caption{Original Block-Diagram in Paper}
  \label{fig:original}
\end{figure}

\section{Open Questions}

\begin{enumerate}
    \item \textbf{Optimization Target of Multi-Domain Traj. Opt.:} 
    What are the optimization targets of Multi-Domain Trajectory Optimization? Do different optimization target types affect the overall performance of the framework?
    \item \textbf{Generalizing to Non-Periodic-Stable Dynamic Motions:}
    Can this framework generalize to non-periodic-stable dynamic motions?
\end{enumerate}

\bibliographystyle{abbrvnat}
\bibliography{main}
\end{document}
